\section{Local kernels and their Frobenius traces}\label{sec:categorical}

We compare compact cohomology of Frobenius-fixed Hecke paths with
dual-group geometry. The comparison has three outputs. At spherical
level, the trace is a pairing with a dg module over the Koszul algebra
for $\operatorname{Ad}(s)e=q^d e$, where $s\in\Gd$ records transport
around the degree-$d$ Frobenius cycle and $e$ is the Lie-algebra
coordinate (Proposition~\ref{cat:loops}). Nef line bundles on the
Springer resolution give the costandard tests used to bound this
pairing (Proposition~\ref{cat:springer}). At Iwahori level, central
kernels and their Wakimoto filtration give Hecke actions and transfers
to spherical level (Proposition~\ref{cat:hecke}). The spherical
description retains dg morphisms; the Iwahori Hecke-algebra action
will be needed on cohomology.

All constructible
categories in this section have their stable, $\mathbf k$-linear enhancements.
Over a field of positive characteristic, $\mathbf k=\overline{\mathbf Q}_\ell$.
Comparisons with complex constructible sheaves are made with characteristic-zero
coefficients. On the dual side put
\[
 S=\operatorname{Sym}(\gd^*[-2]).
\]
Thus a linear polynomial has cohomological degree $2$. Tensor products of
dg modules are derived unless explicitly stated otherwise.

\subsection{Conventions and change of characteristic}

Fix a split maximal torus, an alcove, and its standard parahorics. For
parahorics $P,Q$, write $\mathscr H_{P,Q}$ for the stable category of
equivariant constructible complexes on $P\backslash LG/Q$ generated by
constant orbit complexes and having finite Schubert support. We use the
usual convolution, with compact direct image, for which the point complex
is the unit. The standard and costandard complexes on an Iwahori orbit
of length $n$ are normalized as
\[
 \Delta_w=(j_w)_!\mathbf k[n](n/2),\qquad
 \nabla_w=(j_w)_*\mathbf k[n](n/2).
\]
Fix a square root $q^{1/2}$. Geometric Frobenius acts on
$\mathbf k(-1)$ by $q$. Satake uses the usual correction of the
commutativity constraint by the component parity, so that its target is
ordinary $\operatorname{Rep}(\Gd)$.

\begin{lemma}[Ordinary specialization]\label{cat:specialization}
For every prime $p\ne\ell$, the geometric characteristic-zero and
$\overline{\mathbf F}_p$ categories $\mathscr H_{P,Q}$ admit simultaneous
enhanced equivalences preserving standard, costandard, and intersection
complexes, convolution, and change of parahoric level. They also preserve
restriction of equivariance from $L^+G$ to the pro-unipotent radical $I^u$
of an Iwahori, including the induced maps of mapping complexes.
These are assertions about ordinary geometric categories; Weil structures
will be specified separately.
\end{lemma}

\begin{proof}
Use the split integral group with the given root datum and the strictly
henselian trait $R=\mathbf Z_{(p)}^{\mathrm{sh}}$. Its closed geometric
point is $\overline{\mathbf F}_p$; its geometric generic point has
characteristic zero. First forget the left parahoric equivariance.
On each relative partial affine flag variety $LG_R/Q_R$, use its
\emph{Iwahori-cell} stratification, whose strata are affine spaces over
$R$, and let $\mathscr C_R$ be the thick category generated by their
relative standard constant complexes, with finite Schubert support.
A parahoric orbit is a union of these cells, and its standard constant
complex belongs to $\mathscr C_R$ by localization. We do not regard
parahoric orbits themselves as affine spaces.

Here are the base-change hypotheses being used. Start the Tate lifts
over $S_0=\operatorname{Spec}\mathbf Z[1/\ell]$, which is an admissible
base in \cite[Notation~2.1]{RicharzScholbach}. In that paper's
Section~2.4 the target $T$ of a base change $T\to S_0$ need only be
Noetherian of finite Krull dimension. Thus its Lemma~2.12 applies to
$T=\operatorname{Spec}R$ and to both geometric fibers. It says that
$*$-extension and $!$-restriction of a stratum commute with these
pullbacks on the Tate lifts. Its proof first treats Iwahori cells and
then obtains the parahoric version by localization. Realize these
comparison maps in ordinary etale complexes. Transitivity gives the
maps for restriction from $R$ to either fiber; closing their
equivalence loci under cones and retracts gives the assertion on
$\mathscr C_R$. The same argument shows that restrictions and
corestrictions to a fine cell lie in the thick category generated by
its constant complex. No motivic $t$-structure or new
Beilinson--Soul\'e assertion over $R$ is needed.

On an affine stratum we have
\[
 R\Gamma(\mathbf A_R^n,\mathbf k)=\mathbf k
 =R\Gamma(\mathbf A_{\bar\eta}^n,\mathbf k)
 =R\Gamma(\mathbf A_{\bar s}^n,\mathbf k).
\]
Consequently fiber restriction identifies mapping complexes between
objects generated by its constant complex. Induct on a finite closed
union of strata containing the supports of $A,B\in\mathscr C_R$.
For a closed stratum $i:Z\hookrightarrow Y$ with complementary open
$j:U\hookrightarrow Y$, the fiber sequence
\[
 R\operatorname{Hom}_Z(i^*A,i^!B)
 \longrightarrow R\operatorname{Hom}_Y(A,B)
 \longrightarrow R\operatorname{Hom}_U(j^*A,j^*B)
\]
and the boundary comparison prove full faithfulness on mapping
complexes. Fiber restriction is essentially surjective on the specified
categories because it contains their standard generators and preserves
cones and retracts. This also proves compatibility with composition:
the comparisons are induced by one enhanced restriction functor.

Here is the equivariant version of this argument. On a bounded support
choose a finite-type smooth quotient $K_R$ of the acting parahoric;
its discarded kernel is split pro-unipotent and has trivial positive
equivariant cohomology. Compatible quotients can be used for inclusions
of parahorics and for $I^u\subset L^+G$
\cite[Lemma~2.18]{RicharzScholbach}. Each $K_R$ is an extension of a
split reductive group by split unipotent groups. The cohomology of
$K_R^a$ commutes with restriction to either geometric fiber: this
follows from the Bruhat decomposition of its reductive quotient,
homotopy invariance for the unipotent factors, and localization and
duality for products of affine spaces and split tori. Compute
equivariant mapping complexes by the bar presentation
\[
 Y_R,\quad K_R\times Y_R,\quad K_R^2\times Y_R,\quad\ldots.
\]
At each simplicial degree use the strata
$K_R^a\times\mathbf A_R^n$ coming from the fine Iwahori cells of
$Y_R$. The underlying coefficients in the bar calculation are pulled
back from $Y_R$. The preceding recollement argument therefore applies,
with the single-stratum calculation replaced by that for $K_R^a$.
It gives an equivalence of the mapping complexes at that degree.
Totalizing these equivalences gives the equivariant comparison. No
exchange of a scalar extension with an infinite totalization is used.
The common relative parahoric standard complexes are equivariant
objects; their fiber restrictions generate the specified equivariant
categories. Full faithfulness on the bar mapping complexes and
closure under cones and retracts therefore give essential
surjectivity on those categories as well.
The maps of bar presentations for subgroup inclusions give the claimed
compatibility with restriction of equivariance.

Convolution is twisted exterior product followed by proper direct
image on bounded supports. Smooth descent for the torsor defining the
twisted product and proper base change therefore commute with fiber
restriction. The comparison for all iterated products uses the same
iterated convolution diagram and respects all associativity
compatibilities. This is also the realized convolution comparison of
\cite[Proposition~3.14 and Lemma~3.15]{RicharzScholbach}.
The argument applies to the smooth proper projections between the
partial affine flags. The stratumwise perverse conditions agree on
the two fibers: the restriction and corestriction degrees agree,
and the relative cell dimensions are the same. The equivalence is
therefore perverse $t$-exact and preserves intermediate extensions,
hence intersection complexes. Wakimoto complexes
are finite convolutions of standard and costandard complexes, so they
are preserved as well.

An algebraic closure of the generic field embeds in $\mathbf C$.
On each finite support and each finite bar term, algebraically closed
extension in characteristic zero is fully faithful on bounded
constructible $\overline{\mathbb Q}_\ell$-complexes
\cite[Lemma~A.1]{JKLMBaseChange}.
The etale--Betti comparison is likewise fully faithful
\cite[Corollary~3.4.3, Proposition~3.5.9, and
Remark~3.5.12]{SunComparison}. These comparisons use the derived
pullback and site functors; their isomorphisms on Hom into every shift
give the mapping-complex equivalences. Full faithfulness suffices here:
the specified categories are generated from the matched standard
complexes by finite cones and retracts. Comparing the same bar
presentations at fixed coefficients before totalization gives the
asserted characteristic-zero comparison.
All constructions above use only $p\ne\ell$.
\end{proof}

\begin{theorem}[Derived Satake with Frobenius]\label{cat:satake}
The stable idempotent-complete tensor category generated by the
normalized spherical Satake complexes is monoidally equivalent to
\[
 \operatorname{Perf}^{\Gd}(S).
\]
The normalized Satake object $\operatorname{Sat}(V)$ corresponds to
$S\otimes V$. With the split Weil structures, Frobenius pullback
corresponds to extension of scalars by
\[
 S\longrightarrow S,\qquad \alpha\longmapsto q\alpha
 \quad(\alpha\in\gd^*[-2]),
\]
and acts trivially on the representation labels. These assertions
include the enhanced monoidal structure and all morphisms between
the indicated objects.
\end{theorem}

\begin{proof}
The geometric input for reductive groups is the monoidal dg equivalence
of \cite[Theorem~6.6.1]{BZSV}, which strengthens the monoidal
triangulated calculation of \cite[Theorem~5]{BF}. We use the
corrected version of the latter paper; its purity and multiplicative
formality construction is in \cite[Sections~6.5--6.6]{BF}.
The Galois-compatible formulation in
\cite[Theorem~6.7.5]{BZSV} uses an arithmetic shearing. We verify below
the split Weil normalization and dilation required here.

The regular Satake Ext algebra is the graded equivariant algebra
$\operatorname{Sym}(\gd^*)$, with its linear generators in degree $2$;
the Ext spaces between the other Satake objects are
\[
 \operatorname{Ext}^{\bullet}
       (\operatorname{Sat}(V),\operatorname{Sat}(W))
   =(S\otimes V^*\otimes W)^{\Gd}.
\]
These identifications respect tensor products and composition.
Lemma~\ref{cat:specialization} carries the complete ordinary diagram
of these calculations to every residue characteristic.

For the Weil calculation we need pointwise purity. Over every finite
field, the intersection complex of a split parahoric Schubert variety
is very pure and Tate with semisimple Frobenius
\cite[Section~1.2 and Theorem~2.2.1]{DCHL}. That theorem is stated for
arbitrary split connected reductive groups and all $p\ne\ell$.
Writing $\operatorname{IC}$ here for that paper's unshifted complex,
our normalization is $\operatorname{IC}[d](d/2)$. It says that a
stalk cohomology group in degree $a$ has Frobenius scalar $q^{a/2}$.
Normalized Verdier self-duality gives the same assertion for point
costalks. On a smooth orbit of dimension $d_O$, the identity
$x^!j^!B=(j^!B)_x[-2d_O](-d_O)$ then gives that scalar also on the
degree-$b$ cohomology of the ordinary fiber of $j^!B$.

The equivariant orbit filtration of a morphism complex has local
terms computed from these restrictions and the cohomology of the
split orbit stabilizer $H$. Its reductive quotient is a split Levi,
and the remaining unipotent factors are cohomologically trivial;
$H^c(BH)$ is Tate with scalar $q^{c/2}$. A local term
\[
 \operatorname{Hom}(H^a(A_x),H^b((j^!B)_x))\otimes H^c(BH)
\]
therefore has degree $b-a+c$ and scalar $q^{(b-a+c)/2}$.
All spectral-sequence differentials vanish by their distinct
weights. This proves purity and the semisimplified scalar on the
Ext groups; it alone does not exclude extensions of equal-weight
Frobenius modules. On a degree-two adjoint generator the scalar is
nevertheless exactly $q$: its semisimplification is $q$, and it
commutes with $\Gd$; each simple adjoint summand occurs once. On the
central summand the generators are the first Chern classes of the
split torus, on which Frobenius is also exactly $q$. This determines
Frobenius on the generated algebra. The corresponding goodness
statement for convolution \cite[Corollary~2.2.3]{DCHL}, together with
the perversity of spherical convolution, makes its normalized
tensor multiplicity spaces Tate of weight zero with trivial
Frobenius. Thus the canonical top-cell Weil structures give the
split Satake tensor identifications used above.

For clarity, formality can be applied to this entire diagram at once.
For a mixed morphism complex $C=\bigoplus_w C_w$, use the
multiplicative intermediate complex
\[
 \bigoplus_w\tau_{\leq w}C_w.
\]
Its inclusion into $C$ and its projection to
$\bigoplus_w H^w(C_w)$ are quasi-isomorphisms. Here $C_w$ denotes
the generalized weight summand; no semisimplicity of an arbitrary
pure Weil complex is being inferred. Composition and tensor product
add weights and degrees, so these quasi-isomorphisms are compatible
with every composition and tensor operation. They are
Frobenius-equivariant. Closing this enhanced diagram under finite
cones and retracts gives the assertion. Only its associative
monoidal structure is required below.
\end{proof}

\subsection{Central kernels and Springer tests}

Let $\Lambda=X_*(T)=X^*(\widehat T)$. We call the chamber in which
$J_a$ is standard the positive translation chamber. The opposite
chamber is the costandard chamber. A line bundle on the dual flag
variety is denoted $\mathcal O(\eta)$ with the convention that this
costandard chamber is its nef chamber.

\begin{theorem}[Central and Wakimoto kernels]\label{cat:central}
There is a tensor functor to the Drinfeld center of the homotopy
category
\[
 Z:\operatorname{Rep}(\Gd)\longrightarrow
       \mathcal Z(\operatorname{Ho}(\mathscr H_{I,I}))
\]
with its split Weil structure. Its interchanges are compatible with
projection to every standard parahoric. At a hyperspecial parahoric
the projected functor is Satake.

There are invertible kernels $J_a$, $a\in\Lambda$, with coherent
multiplication $J_a\star J_b\simeq J_{a+b}$. They are standard on
the positive chamber and costandard on its opposite. The central
object $Z(V)$ has a finite functorial Wakimoto filtration, tensor
compatible in $V$, with
\[
 \operatorname{gr}_a Z(V)=J_a\otimes V_a.
\]
Its induced central interchange with $J_b$ is the ordinary interchange
of the two translation factors and the identity on the multiplicity
space, with the fixed Satake parity convention. The multiplicity
spaces have trivial semisimplified Frobenius. These statements hold
over every $\mathbf F_q$ under consideration.
For a fixed central label, its interchange in the other kernel
variable is a chosen natural equivalence of exact enhanced functors.
We use the central tensor identities in the homotopy category;
we do not assert an enhancement of $Z$ to an infinity-categorical
Drinfeld center.
\end{theorem}

\begin{proof}
The central construction and the Wakimoto filtration for every split
reductive group over the base in question follow from
\cite[Theorem~1.4 and Theorem~5.24]{CvdHS}, by realization.
The scope of its centrality assertion is exactly the homotopy
category. More precisely, its Theorem~4.38 constructs the two
natural maps in its Equation~(4.7) from the enhanced nearby-cycles
and convolution functors, using Lemma~4.20 and Theorem~4.36(3).
For a fixed central label these give the stated natural equivalence
of exact functors in the other kernel variable, so it applies to
an actual cofiber diagram. Theorem~4.38 supplies the homotopy
hexagon, and Theorem~4.41 the homotopy tensor compatibility.
These finite identities are the ones needed for products of
character operators on cohomology below. The maps are defined
over $\mathbf F_q$ and carry their actual Weil structures, after
fixing the nearby-cycles Frobenius lift.

We record explicitly the graded-interchange argument, to include all
root data. The Wakimoto graded category is the category of
$\Lambda$-graded vector spaces, and the graded central functor has
the weights of restriction from $\Gd$ to $\widehat T$. For every
dominant character $\lambda$ of the actual torus there is its
highest-weight quotient to the line labelled by $J_\lambda$.
On this quotient the central interchange is the ordinary
interchange: restrict the fusion diagram to the open Cartan-position
cell, where the sheaves are normalized constants and the quotient
is the highest-line restriction. The two tensor generators are
interchanged there, with exactly the fixed Satake parity correction.
Naturality and the central tensor identity now imply the same
interchange for an arbitrary central label with $J_\lambda$.
Indeed the two opposite central interchanges have product one,
and passage to this quotient identifies one of them with the
ordinary interchange, so identifies the other with its inverse.
This is the argument of \cite[Lemma~18 and Section~3.6.6,
Equation~(21)]{AB}. Dominant characters generate $\Lambda$ as a
group. Tensor compatibility and invertibility of the $J_\lambda$
extend the assertion to every $\lambda$. The argument uses the
actual character lattice, and applies to products and central torus
factors as stated. Its top-cell maps are defined over $\mathbf F_q$,
so it preserves the specified normalization and Weil maps.
Projection compatibility follows from proper base change for the
nearby-cycles construction: the projection of flag levels is a
smooth proper flag fibration, and the commutation before taking
nearby cycles is the commutation of modifications at disjoint points.
Using the same multiple-point diagram makes these compatibilities
simultaneous for all the parahorics and all products.

The split Weil multiplicities can also be checked without any
choice of splitting of the filtration. The trace of the central
object is the Bernstein central character, over each finite
extension. The Wakimoto trace functions are linearly independent.
Consequently all power traces on $V_a$ equal $\dim V_a$;
their semisimplified eigenvalues are one. This statement permits
unipotent extensions, which do not affect the eigenvalue
separations used later.
\end{proof}

\begin{proposition}[The spherical Springer test]\label{cat:springer}
Let $\operatorname{For}$ forget $L^+G$-equivariance to
$I^u$-equivariance on $\operatorname{Gr}_G$, without changing the
underlying complex, and let $\operatorname{Av}_*$ be its right
adjoint. Put $W_\eta=J_\eta\star\delta_0$. For every integral
weight $\eta$ in the closed costandard chamber, derived Satake
identifies
\[
 \operatorname{Av}_*W_\eta
 \quad\text{with}\quad
 \bigl(R\pi_*\mathcal O_{\widetilde{\mathcal N}}(\eta)\bigr)_{\rm sh},
 \qquad
 \pi:\widetilde{\mathcal N}\longrightarrow\gd.
\]
The subscript means that coherent degree $a$ and polynomial degree
$b$ are assigned total degree $a+2b$. This is an equivalence of
ordinary equivariant dg $S$-modules. The $S$-action is polynomial
restriction along $\pi$, and the assertion is compatible with
tensoring by representations and with all morphisms between free
Satake labels. The displayed convention introduces no extra shift.
Any smooth-pullback normalization of $\operatorname{For}$ changes
the formula only by a fixed shift and Tate twist, independent of
$\eta$.
\end{proposition}

\begin{proof}
We give the algebra-to-module calculation, including its root-datum
and equivariance conventions. Let $\mathcal R$ be the regular Satake
ind-object. Its second regular action gives the $\Gd$-action on the
following algebras and modules. The regular calculations are
\[
 \operatorname{Ext}^{\bullet}_{L^+G}(\delta_0,\mathcal R)=S,
 \qquad
 \operatorname{Ext}^{\bullet}_{I^u}(\delta_0,\mathcal R)
       =\mathbf k[\mathcal N]_{\rm sh},
\]
and the Wakimoto calculation for nef weights, including its products, is
\begin{equation}\label{cat:section-algebra}
 \bigoplus_{\eta\ \text{nef}}
 \operatorname{Ext}^{\bullet}_{I^u}
            (\delta_0,J_\eta\star\mathcal R)
 \simeq
 \left(\bigoplus_{\eta\ \text{nef}}
       \Gamma(\widetilde{\mathcal N},\mathcal O(\eta))\right)_{\rm sh}.
\end{equation}
The algebra calculation is
\cite[Theorems~7.3.1 and 7.6.1]{ABG}; the section algebra and its
product maps are \cite[Theorem~8.5.2 and Equations~(8.7.2)--(8.7.8)]{ABG}.
We explain why using the actual lattice $\Lambda$ entails no
isogeny restriction here. Wakimoto inversion identifies the left
side for $\eta$ with the appropriate fixed-point costalk of
$\mathcal R$. Its cohomological filtration is the principal-nilpotent
filtration on the $-\eta$ weight space of $\mathbf k[\Gd]$.
The fixed-point and filtration assertions of
\cite[Propositions~5.5.2 and 5.6.2]{GinzburgLoop} are stated for the
full cocharacter lattice of an arbitrary semisimple group.
The section-filtration identification used in
\cite[Equations~(8.6.2)--(8.6.3)]{ABG} applies to every dominant
character of the actual dual torus. Multiplication of the regular representation,
followed by the fixed-point-to-convolution base-change maps in
\cite[Equations~(8.7.6)--(8.7.8)]{ABG}, identifies the products in
Equation~\eqref{cat:section-algebra}. This proves the required
calculation for every semisimple isogeny type, without an assertion
about all objects in a generalized ABG equivalence.

For a torus the calculation is immediate: its Grassmannian is
$\Lambda$, its pro-unipotent group is cohomologically trivial, and
the algebra map is $\operatorname{Sym}(\operatorname{Lie}(\widehat
T)^*[-2])\to\mathbf k$. For a general reductive group, perform
the characteristic-zero calculation using its central isogeny to
the product of its adjoint group and its abelianization. Each
Grassmannian component identifies with the corresponding component
in that product; finite central kernels have trivial positive
cohomology with these coefficients. On the dual side this restricts
the regular representation and its products to the prescribed
central-character lattice. It therefore restricts precisely the
preceding section-algebra calculation. Central translations have
length zero, and central linear variables map to zero. Finally
Lemma~\ref{cat:specialization} transports this ordinary enhanced
diagram, including the forgetful map of algebras and the module
action, to every $p\ne\ell$. The regular ind-object causes no
problem: mapping out of the point in these Ind-completions commutes
with its filtered union.

We now identify the polynomial map. The equivariant-cohomology
spectral sequence for the regular internal Hom has second page
\[
 H^{\bullet}(BG)\otimes
       \operatorname{Ext}^{\bullet}_{I^u}(\delta_0,\mathcal R).
\]
Both factors are even, so its edge map induces
\[
 \mathbf k\otimes_{H^{\bullet}(BG)}
       \operatorname{Ext}^{\bullet}_{L^+G}(\delta_0,\mathcal R)
 \simeq
       \operatorname{Ext}^{\bullet}_{I^u}(\delta_0,\mathcal R).
\]
Freeness over $H^{\bullet}(BG)$ follows by lifting a homogeneous
basis on this second factor and comparing its associated graded.
Taking $\Gd$-invariants identifies $H^{\bullet}(BG)$ with
$S^{\Gd}$, since the invariant part of $\mathcal R$ is the unit.
The map thus has kernel $S(S^{\Gd})_{>0}$ and is restriction to
the nilpotent cone. On degree-two simple adjoint factors its
identification differs from the usual one at most by nonzero
scalars, by equivariance and surjectivity. Absorb these constant
scalars in the choice of coordinates; they commute with the
Frobenius dilation. This also identifies the action on each module
in Equation~\eqref{cat:section-algebra}.

The comparisons just used are comparisons of dg modules. In
characteristic zero the simultaneous purity construction of
\cite[Proposition~9.5.2 and Lemma~9.5.3]{ABG}, or the multiplicative
weight truncation in Theorem~\ref{cat:satake}, formalizes the
algebra, its map, and its positive modules together. The ordinary
enhanced specialization carries that diagram to the desired
characteristic. No identification of an arbitrary Weil structure
on a Springer test is needed for the present assertion.

For every finite-dimensional $V$, adjunction now gives
\begin{align*}
 R\operatorname{Hom}(\operatorname{Sat}(V),\operatorname{Av}_*W_\eta)
 &\simeq R\operatorname{Hom}(\operatorname{For}\operatorname{Sat}(V),W_\eta)\\
 &\simeq
   \bigl(R\Gamma(\widetilde{\mathcal N},
                    V^*\otimes\mathcal O(\eta))\bigr)^{\Gd}_{\rm sh}\\
 &\simeq R\operatorname{Hom}_{S,\Gd}
       (S\otimes V,(R\pi_*\mathcal O(\eta))_{\rm sh}).
\end{align*}
For every nef $\eta$, including walls, the higher line-bundle
cohomology vanishes by \cite[Theorem~2.4 and Remark~2.7(2)]{Broer}.
Thus the section-module computation indeed computes the derived
pushforward appearing here. This use of vanishing is entirely on
the characteristic-zero dual side. The polynomial computation
proves naturality in all morphisms between the free labels, which
compactly generate. Yoneda proves the asserted equivalence.
The pushforward is coherent on the smooth affine space $\gd$;
an equivariant graded finite free resolution, followed by shearing,
shows that it is a perfect $S$-module.

Finally our forgetful functor preserves the underlying sheaf. If
one instead uses $\operatorname{For}[c](c/2)$, its right adjoint
is $\operatorname{Av}_*[-c](-c/2)$. The integer $c$ is the
dimension fixed by that change of level, independently of $\eta$.
\end{proof}

\subsection{The geometric cyclic functor}

Bundle stacks below may have fixed full congruence level away from
the selected legs and fixed parahoric level at them. They may also
be divided by a discrete lattice of central modifications, as in
Section~\ref{sec:rotation}. At a closed place every geometric leg
is retained, with Frobenius cyclically permuting them. A kernel
may change these levels. It always has bounded relative position.

For a composite kernel $L$ on a bundle stack $\mathcal B$, define
\[
 \mathsf K(L)=R\Gamma_c(\mathcal B,
                     \operatorname{Graph}(\mathrm{Fr})^*L).
\]
Equivalently this is compact cohomology of the stack of paths from
$P$ to ${}^{\mathrm{Fr}} P$, with all intermediate bundles and all Hecke
coefficients retained. Fixed kernels at other places may be
included in every path. Compact cohomology of a locally finite
type stack is the colimit over finite-type open substacks, using
extension by zero.

\begin{proposition}[Coherent path rotation]\label{cat:cyclic}
The preceding construction is an exact enhanced Frobenius-twisted
cyclic functor. Put $\Phi=\mathrm{Fr}^*$ on local kernels. With convolution
ordered in the direction of the path, its cyclic maps are
\begin{equation}\label{cat:last-first}
 \rho_{U,V}:\mathsf K(U\star V)
       \xrightarrow{\sim}\mathsf K(\Phi(V)\star U).
\end{equation}
They are natural in all enhanced morphisms and coherently
compatible with units, iterated products, and change of level.
Here $\rho_{U,\mathbf1}$ is the identity. The full slide
$\rho_{\mathbf1,V}:\mathsf K(V)\to\mathsf K(\Phi V)$,
followed by the Weil structure of $V$, is cohomological
Frobenius. Slides at disjoint sets of legs commute; their product
is the full slide.
\end{proposition}

\begin{proof}
For two kernels, the path description is
\[
 (P,Q;\ U(P,Q),\ V(Q,{}^{\mathrm{Fr}}P)).
\]
Rotating the first step sends $(P,Q)$ to $(Q,{}^{\mathrm{Fr}}P)$.
To obtain pullback in that direction use its inverse, which sends
$(P,Q)$ to $({}^{\mathrm{Fr}^{-1}}Q,P)$. The two coefficients on this
rotated path are $\Phi(V)({}^{\mathrm{Fr}^{-1}}Q,P)$ and $U(P,Q)$.
This is precisely Equation~\eqref{cat:last-first}. In this
description $\mathrm{Fr}^{-1}$ may be interpreted on perfections. Etale
constructible sheaves and their compact direct images are
invariant under the resulting universal homeomorphisms, so the
construction also defines the same map before perfection.

For any finite list of kernels use the stack of all its
intermediate bundles. Multiplication forgets an intermediate
bundle, a unit repeats one, and rotation moves an end bundle
through Frobenius. The identities between these operations are
identities of the corresponding path diagrams. Applying compact
direct image, base change, and the projection formula in the
enhanced six-operation formalism preserves these identities
coherently. The required precise finite-type foundations are
the enhanced adic operations and their base-change and projection
formula compatibilities in
\cite[Section~7.1 and Theorem~7.2.7]{LiuZheng}.

One can express the coherence as a balanced bar construction.
In simplicial degree $n$ retain all intermediate kernels in the
path; inner faces compose adjacent kernels, degeneracies insert
units, and the end face is the rotation just described.
Their common refinements are again the same path stack, which
supplies every simplicial identity and its higher compatibility.
This verifies coherence in the enhancement, rather than only
pairwise isomorphisms in its homotopy category.

The same reasoning applies to the locally finite type stack.
A fixed finite diagram uses finitely many bounded kernels.
The image and inverse image of a finite-type bounded open under
the relevant bounded-position maps are contained in some larger
finite-type bounded open. Thus that diagram is computed on a
finite-type stage, after enlarging stages when necessary.
The finite-type maps are compatible under extension by zero,
so their filtered colimit preserves the same identities.
Central-lattice quotients are treated by the corresponding
locally finite sums of these diagrams.

When $U=\mathbf1$, inverse rotation of the path is $\mathrm{Fr}^{-1}$.
Its compact direct image is pullback by $\mathrm{Fr}$. Composing with the
Weil structure is therefore exactly cohomological Frobenius.
Moving disjoint groups of labels gives commuting rotations of
the common path diagram, with product the complete rotation.
\end{proof}

\subsection{Hecke characters and finite flag transfers}

If $U$ is a rigid Weil kernel with left dual $U^\vee$ and with
an interchange $U\star L\simeq L\star U$, define its character
operator on $\mathsf K(L)$ by
\begin{equation}\label{cat:character}
\begin{split}
 \mathsf K(L)&\longrightarrow
 \mathsf K(L\star U^\vee\star U)
 \xrightarrow{\rho}
 \mathsf K(\Phi U\star L\star U^\vee)\\
 &\longrightarrow\mathsf K(U\star L\star U^\vee)
 \longrightarrow\mathsf K(L\star U\star U^\vee)
 \longrightarrow\mathsf K(L).
\end{split}
\end{equation}
The arrows use coevaluation, rotation, the Weil structure,
interchange, and evaluation, respectively. We denote this
operator by $\chi_U$. The same formula for a rigid kernel
between two different levels is called transfer.
The individual formula is a morphism of complexes. The Iwahori
Hecke-algebra action asserted next is its action on
$H^\bullet\mathsf K(L)$; no enhanced action of that algebra is
deduced from the homotopy central functor.

\begin{proposition}[Hecke action and transfer]\label{cat:hecke}
Suppose that $L$ has a central label at a closed place $x$.
The operators in Equation~\eqref{cat:character} give the usual
affine Hecke algebra at $x$, with parameter $q_x$, on
$H^\bullet\mathsf K(L)$. They give ordinary Hecke operators on
$H^\bullet\mathsf K(\mathbf1)$.

More explicitly, let $H_w$ denote the operator with characteristic
function $1_{IwI}$ and Iwahori volume one. Then
\[
 H_s^2=(q_x-1)H_s+q_x,
 \qquad H_wH_{w'}=H_{ww'}\quad
       \text{if }\ell(ww')=\ell(w)+\ell(w').
\]
The positive normalized Bernstein translations are
\begin{equation}\label{cat:translation-sign}
 T_a=(-1)^{\ell(t^a)}\chi_{J_a},
 \qquad
 T_a=q_x^{-\ell(t^a)/2}H_{t^a}\quad(a\text{ positive}),
\end{equation}
and extend multiplicatively to the entire translation lattice.

For a standard parahoric $P\supset I$, its two level transfers
identify the hyperspecial or parahoric trace cohomology with the image of
\[
 e_P=\left(\sum_{w\in W_P}q_x^{\ell(w)}\right)^{-1}
                  \sum_{w\in W_P}H_w.
\]
At hyperspecial level the spherical character action agrees
with the Bernstein center. The individual-leg slides commute
with these actions and transfers. On the Wakimoto graded terms
the Bernstein action is the translation action tensored with
the weight multiplicities in Theorem~\ref{cat:central}.
The Wakimoto filtration gives long exact sequences equivariant for
these Bernstein translations and for the selected-leg slides.
Its graded terms carry the stated translation action. Thus flat
completion, ordered-weight projections, and spectral projectors
may be applied to these cohomology sequences.

There is also the following geometric form of the rank-one
assertion. If an exact translation position is parallel to the
simple wall of $s$, the operator $H_s+1$ on its path trace is
pullback and sum along the finite etale family of Frobenius-fixed
rank-one refinement flags, of degree $q_x+1$. Consequently
$H_s$ takes the sum over the other flags.
\end{proposition}

\begin{proof}
The kernels in use are rigid. Standard simple kernels are
invertible with inverse the corresponding costandard kernel;
length-zero kernels are invertible as well. Finite standard
filtrations, cones, and retracts preserve dualizability. The
projection kernels have the pullback and proper-direct-image
adjunctions of the smooth proper flag projection.

The character formula is multiplicative under convolution:
the coevaluation for a product is the composite of its two
coevaluations, and the cyclic and interchange identities move
the two factors successively. Applying the triangle identities
then gives the product of their characters on cohomology. This
uses only the homotopy hexagon and tensor identities specified
in Theorem~\ref{cat:central}. For additivity, fix the central
label $L$ and apply its chosen natural equivalence of exact
enhanced functors to the cofiber diagram of the test kernels.
This supplies a morphism of that cofiber diagram, not merely
unrelated maps on its three objects. On a mapping-cone model
the interchange is triangular with respect to the two summands, while
coevaluation and evaluation pair equal summands. The two
diagonal contributions are the character of the second term
and the negative of the character of the first term. This
also explains why shifts give their usual alternating signs.

For the quadratic relation use the unshifted closed kernel
$C_s=\mathbf k_{P_s/I}$, where $P_s/I\simeq\mathbf P^1$.
Convolution through the partial level gives
\[
 C_s\star C_s\simeq
 R\Gamma(\mathbf P^1,\mathbf k)\otimes C_s.
\]
The open--closed triangle gives $\chi_{C_s}=H_s+1$.
The Weil trace of $R\Gamma(\mathbf P^1,\mathbf k)$ at $x$
is $1+q_x$. Thus $(H_s+1)^2=(1+q_x)(H_s+1)$, the asserted
relation. Length-additive convolution of open cells is an
isomorphism, giving the length and braid relations. For a
degree-$d_x$ place, all kernels are external products on its
geometric legs with cyclic Frobenius. The trace of a graded
cycle on $C^{\otimes d_x}$ is the trace of the composite
Frobenius on $C$; hence the parameter is $q^{d_x}$.
The same identity shows that a shifted simple or translation
kernel contributes its single-factor sign
$(-1)^{\ell(t^a)}$, not a spurious sign depending on the
number of geometric legs. This proves
Equation~\eqref{cat:translation-sign}. Length parity is a
character of the translation lattice, so multiplicativity is
preserved.

At $L=\mathbf1$ the path stack is the Frobenius-fixed bundle
groupoid. On it, the unit, counit, and Weil structure in
Equation~\eqref{cat:character} sum the stalk traces over the
Hecke correspondence, with its stack multiplicities. This is
the ordinary action on compactly supported functions on the
rational bundle groupoid. The construction for other central
labels uses the same diagrams and their central interchange,
so the verified relations apply there as well.

For general $P$ repeat the closed-kernel computation with the
full flag $P/I$. Connected transition groups act trivially
on its cohomology, so its cohomology is the constant Tate
complex with Frobenius trace
$c_P=\sum_{w\in W_P}q_x^{\ell(w)}$. The composite of the
two transfers at the $P$-level is multiplication by $c_P$;
the opposite composite is $\sum_{w\in W_P}H_w$.
Since $c_P\ne0$, the two normalized transfers give the
stated idempotent identification. Theorem~\ref{cat:central}
identifies the projected central kernel with Satake at a
hyperspecial level. Trace additivity on its Wakimoto
filtration then gives the usual Bernstein central character,
which proves compatibility of the spherical action.

We spell out the assertion about parallel walls, which is
needed for a geometric correspondence later. For a translation
$t^a$ parallel to the $s$ wall, its character on the rank-one
root subgroup has valuation zero. Thus the exact partial
relative position identifies the rank-one refinement flag at
the two ends. On the Frobenius path stack the allowed flags
are the fixed flags of that semilinear identification. They
form a finite etale family $Z\to Y$ of degree $q_x+1$.
The closed kernel factors through the partial level, and
its character factors through the two transfer maps.
Those maps are the pullback and sum for $Z\to Y$.
Indeed, smooth proper base change reduces the unit/counit
calculation to the diagonal of $\mathbf P^1$ and the graph
of its twisted Frobenius. Their intersections are transverse:
Frobenius has zero tangent map and the diagonal contributes
the identity. Every fixed flag therefore has coefficient
one. Etale locally on $Y$, where $Z$ is a disjoint union of
sections, the resulting maps are the diagonal map and the
sum map between constant coefficient complexes. They agree
on overlaps and hence descend. This calculation is unchanged
after tensoring with the coefficient pulled back from the
partial path stack, or after retaining automorphisms of the
bundles. Subtracting the diagonal, by the open--closed
triangle, gives the other-flags operator $H_s$.

Finally, naturality of full rotation makes Frobenius commute
with every Weil Hecke action. A slide at an auxiliary disjoint
place commutes with this action by its common disjoint-leg
diagram. Dividing the full slide by the other, invertible
slides proves the same assertion for the selected-leg slide.
Projection transfers use those same diagrams. The graded
translation assertion follows from the map compatibility in
Theorem~\ref{cat:central}. More explicitly, those filtered
interchanges and rotation maps commute with the inclusion,
quotient, and connecting maps of each successive filtration
triangle. Applying cohomology gives the asserted equivariant
long exact sequences. This is all the filtration compatibility
needed here; it does not require an action on a derived
completion of the Iwahori trace.
\end{proof}

Whenever an unqualified equality between a translation slide
and Frobenius is used below, the large translation labels are
chosen divisible by an even integer. Their length parity then
vanishes. The weights of a representation with such an
extremal label have the same even parity, since their
differences are roots of the dual group. Filter translations
remain arbitrary and always use the explicit normalization
in Equation~\eqref{cat:translation-sign}.

\subsection{The algebra of a twisted loop}

The final local calculation is given in terms of dg algebras,
so positive cohomological polynomial degrees cause no
connectivity assumption. It applies to the entire compact
cohomology functor, including fixed disjoint kernels.
Its enhanced input is derived spherical Satake and the coherent
path functor, independently of the homotopy central functor at
Iwahori level.

\begin{proposition}[Spectral form of a Frobenius trace]\label{cat:loops}
At a hyperspecial closed-place leg of degree $d$, the functor
$\mathsf K$ on perfect spherical labels has the form
\begin{equation}\label{cat:loop-pairing}
 \mathsf K(P)\simeq
 \left(\mathcal A\otimes_{B_{q^d}}
                  (B_{q^d}\otimes_{S^{\otimes d}}P)\right)^{\Gd},
\end{equation}
where $\mathcal A$ is an equivariant dg module, with no
coherence or boundedness assumption. After eliminating
intermediate factors, the algebra is
\begin{equation}\label{cat:loop-algebra}
 B_r=\mathcal O(\Gd)\otimes
       \operatorname{Sym}(\gd^*[-2]\oplus\gd^*[-1]),
 \qquad r=q^d.
\end{equation}
Write $s$ for its degree-zero holonomy variable and $e$ for
its degree-two linear variable. The second summand supplies
odd Koszul generators $\eta_\alpha$ of degree $1$, with
\begin{equation}\label{cat:holonomy}
 d\eta_\alpha=
       \langle\alpha,\operatorname{Ad}(s)e-q^d e\rangle.
\end{equation}
The group acts by simultaneous conjugation.

The $d$-fold selected-leg slide on representation labels is
the action of $s$ on each factor. Spherical Hecke characters
act by the corresponding invariant characters of $s$.
On cohomology these descriptions agree with the transfers in
Proposition~\ref{cat:hecke}. A degree-two morphism linear
in the selected $e$ variable multiplies the corresponding
slide eigenvalue by $q^d$. All statements are natural in
perfect labels and their dg morphisms.
\end{proposition}

\begin{proof}
Put $\mathcal C=\operatorname{Mod}_S
(D(\operatorname{Rep}_{\rm rat}(\Gd)))$ and let $\Phi_r$
be extension of scalars by $\alpha\mapsto r\alpha$.
The free objects $S\otimes V$ are compact, dualizable
generators: rational invariants of a reductive group in
characteristic zero are exact and commute with direct sums
\cite[Remark~15.13 and Proposition~15.16]{MilneAlgebraicGroups}.
Thus this category is the Ind-completion of its perfect
objects.

We first calculate the universal trace for one factor.
Let $\mathcal C^e=\mathcal C\otimes\mathcal C^{\rm rev}$.
Use the regular right action
$L\cdot(P,Q)=Q\otimes L\otimes P$ and the twisted left
action $(P,Q)\cdot M=\Phi_r(P)\otimes M\otimes Q$.
The pairing $(L,M)\mapsto\mathsf K(M\otimes L)$ is
coherently balanced: the balancing map is precisely
\[
 \mathsf K(M\otimes Q\otimes L\otimes P)
 \xrightarrow{\rho}
 \mathsf K(\Phi_r(P)\otimes M\otimes Q\otimes L).
\]
Proposition~\ref{cat:cyclic} verifies all the bar identities.
Consequently $\mathsf K$ extends to a continuous functional
on the relative tensor product of these two module
categories.

Here is an explicit algebra calculation of that relative
tensor product. It is the derived intersection of the
diagonal and the graph $(\Phi_r,\operatorname{id})$ in the
product presentation of $[\operatorname{Spec}S/\Gd]$.
This language abbreviates a module-algebra computation:
\cite[Proposition~4.1]{BFN} identifies the tensor product of
module categories with modules over the derived tensor
product of their defining algebra objects. It applies in
the stable presentable category of equivariant modules,
without a connectivity condition.

An object of this intersection has a diagonal coordinate
$x$, a graph coordinate $y$, and two group identifications
\[
 a:x\longrightarrow r y,
 \qquad b:x\longrightarrow y.
\]
Resolve the equations
$\operatorname{Ad}(a)x-r y=0$ and
$\operatorname{Ad}(b)x-y=0$ by their degree-one Koszul
generators. Their algebra, before group descent, is the
polynomial algebra on $x,y$ and the two group variables
with those Koszul generators. These resolutions are
$K$-flat: their exterior-degree filtrations have free
graded polynomial quotients. Eliminate $y$ and its
Koszul generator using the second equation. This is an
invertible linear substitution followed by removal of
contractible Koszul pairs. Trivialize $b$ by the product
group action, put $e=y$ and $s=a b^{-1}$, and transport
the remaining generator by the same change of basis.
The residual group is the diagonal $\Gd$, its action is
conjugation, and the remaining differential is
\[
 d\eta_\alpha=
       \langle\alpha,\operatorname{Ad}(s)e-r e\rangle.
\]
This is exactly $B_r$. In particular the specified
null-homotopies of the equation have been retained.

The universal trace sends a label $P$ to $B_r\otimes_S P$.
The comparison from the second-coordinate representation
to the first-coordinate representation in the preceding
diagram is $a b^{-1}:y\to r y$. It is the last-to-first
slide of Equation~\eqref{cat:last-first}. Hence its action
on a representation is $V(s)$, with the asserted
orientation. Coevaluation and evaluation then give
$\operatorname{Tr}_V(s)$ for a rigid character transfer.

For a cyclic degree-$d$ leg, retain $d$ coordinates and
the $d$ twisted identifications. Trivializing all but one
of the group variables eliminates the intermediate
coordinates and their Koszul pairs. The remaining
transport is their ordered product $s$, and its equation
is $\operatorname{Ad}(s)e=q^d e$. Carrying a representation
label around this full cycle acts by $V(s)$. For an
external tensor label every factor has this same full
transport; a single partial turn also performs the
appropriate graded cyclic permutation. This explains
both the $d$-fold slide and the map from the original
$S^{\otimes d}$ to the loop algebra in
Equation~\eqref{cat:loop-pairing}.

It remains to justify that the continuous functional is
pairing with a module, including when it is not coherent.
Let $\mathcal T=\operatorname{Mod}_{B_r}
(D(\operatorname{Rep}_{\rm rat}(\Gd)))$ and write the
functional as $L:\mathcal T\to D(\mathbf k)$.
Again $B_r\otimes V$ are compact dualizable generators.
The exact contravariant functor on $\mathcal T^c$
given by $P\mapsto L(P^\vee)$ is represented by an
object $\mathcal A\in\operatorname{Ind}(\mathcal T^c)
=\mathcal T$. Tensor duality gives, for compact $M$,
\[
 L(M)\simeq
 \operatorname{Hom}_{\mathcal T}(M^\vee,\mathcal A)
 \simeq(\mathcal A\otimes_{B_r}M)^{\Gd}.
\]
Both sides preserve colimits, proving the formula for
all $M$, and hence Equation~\eqref{cat:loop-pairing}.

The character calculation is natural in labels and
their evaluation and coevaluation maps. The projection
transfer of Proposition~\ref{cat:hecke} is therefore the
same character operation on cohomology in this model. Finally, for a
linear degree-two morphism $u:P\to Q[2]$, cyclic
naturality and Theorem~\ref{cat:satake} give
\[
 \sigma_Q[2]\,\mathsf K(u)
       =q^d\,\mathsf K(u)\,\sigma_P,
\]
where $\sigma$ is the full selected-leg slide. This
proves the stated change in eigenvalue. The degree is
even, so no additional Koszul sign occurs.
\end{proof}

In particular $Z_{\rm sph}=\mathcal O(\Gd)^{\Gd}$ acts as an
actual degree-zero central dg algebra through $B_r$. Here is the
precise spectral-projector consequence. For an external Satake label
with representation $V^{\otimes d}$, let
$\rho=V^{\otimes d}$ and let $\sigma$ be its $d$-fold selected-leg
slide. Before tensoring with $\mathcal A$ and taking invariants,
$\sigma$ is the matrix $\rho(s)$ on the finite free $B_r$-module
$B_r\otimes V^{\otimes d}$. For any positive integer $N$,
\[
 p(z)=\det\bigl(z-\rho(s)^N\bigr)\in Z_{\rm sph}[z]
\]
annihilates $\sigma^N$ strictly by Cayley--Hamilton. The identity
persists after tensoring with $\mathcal A$ and taking invariants.
If a flat $Z_{\rm sph}$-algebra $R$ gives a coprime factorization
$p=p_0p_1$, choose $a p_0+b p_1=1$. Then
$b(\sigma^N)p_1(\sigma^N)$ is an actual chain idempotent on the
spherical complex after this base change. It selects the $p_0$
factor and gives its dg direct summand, regardless of boundedness.
This is the enhanced projector used after transferring a
cohomological Iwahori summand to hyperspecial level.

Two consequences will be used repeatedly. First, imposing
finitely many equations in invariant functions of $s$ is
ordinary derived tensor with a finite Koszul complex and
commutes with every operation just described. Second,
after such a specialization, an invertible coefficient
of a nonresonant linear equation permits its variable
and Koszul generator to be eliminated as a contractible
pair. Thus the action of the original polynomial
variables on a resonant slice is their ordinary
restriction, with cohomological degree still $2$.
